The ten, its mode vectors, the number operators and the Majorana element $T = v_{1}v_{2}v_{3}v_{4}v_{5}$ are those of Proposition \claimref{standard-model.majorana-element}; $T_{3R}$, $B - L$ and $Y = T_{3R} + \tfrac12(B - L)$ are as in Proposition \claimref{synthesis.georgi-jarlskog-is-rung-vector-times-pair-part-of-majorana-square} and Proposition \claimref{standard-model.majorana-element}(2), (4) (there $X = 2T_{3R} - \tfrac32(B - L) = 2Y - \tfrac52(B - L)$, which is this relation). The frame boost, the lift family and the $*$-norm of the rung component are those of Theorems \claimref{synthesis.rung-unit-is-frame-boost}(4) and \claimref{synthesis.rung-unit-transverse-to-norm}(1). The stored boost acts on the rung plane of slot $4$; the setting takes its $U(2)$-covariant extension to both rung slots, the identity on the pair slots, and (5) treats the two other extensions. The same boost is read as the conformal dilatation of the ground point in the result \claimref{synthesis.scale-free-rung-potentials-need-the-flat-physical-branch}(5); on the rung plane the two readings coincide.

(1) Since $Q(v_{a}) = Q(w_{a}) = 0$ and $v_{a}w_{a} + w_{a}v_{a}$ is unchanged when $v_{a} \mapsto sv_{a}$, $w_{a} \mapsto s^{-1}w_{a}$, the map $B_{s}$ on vectors is an isometry, so it extends uniquely to an automorphism of the Clifford algebra; the check verifies that the images of the ten generators obey the Clifford relations for $s = 2, \tfrac32$. On the lift family, $c_{1} = -w_{4} + \tfrac{\Lambda}{3}v_{4}$ is sent by $B_{1/\lambda}$ to $-\lambda w_{4} + \tfrac{\Lambda}{3}\lambda^{-1}v_{4} = c_{\lambda}$. From the anticommutators, $[n_{a}, v_{b}] = \delta_{ab}v_{b}$ and $[n_{a}, w_{b}] = -\delta_{ab}w_{b}$, so on a product of mode vectors $\operatorname{ad} N_{R}$ has eigenvalue $q$, the number of rung $v$ factors minus the number of rung $w$ factors, and $B_{s}$ multiplies that product by $s^{q}$. Every element is a sum of such products, so $B_{s}$ acts on the $\operatorname{ad} N_{R}$ eigenspace of weight $q$ as $s^{q}$. Since $N_{R} = 2T_{3R} + 1$, $\operatorname{ad} N_{R} = 2\,\operatorname{ad} T_{3R}$, and $B_{e^{\eta}} = \exp(\eta\,\operatorname{ad} 2T_{3R})$: the frame boost is the flow of the right-handed isospin at real rapidity $2\eta = t$. (The compact flow $\exp(I_{4}\varphi\,\operatorname{ad} 2T_{3R})$ is the unitary one, numbers and the module's complex structure are read in $\mathrm{span}\{1, I_{4}\}$, the module read on the tower's floors, where $I_{4}$ commutes with the tower vectors; by Theorem \claimref{synthesis.rung-unit-is-frame-boost}(1) the hypercharge already acts on the real rung summand as a boost.) The check verifies the eigenvalue relation and $B_{s}X = s^{q}X$ on nine monomials, including $T$ and $T^{*}$.

(2) On a product of mode vectors let $p$ be the number of pair $v$ factors minus pair $w$ factors and $r$ the same for the rung slots. Then the weights are additive: $\operatorname{ad}(B - L)$ has weight $-\tfrac23 p$, $\operatorname{ad} T_{3R}$ has $\tfrac12 r$, $\operatorname{ad} N_{R}$ has $r$ and $\operatorname{ad} Y$ has $\tfrac12 r - \tfrac13 p$; the check verifies these four weights by direct commutators on all $3^{5}$ slot types (each slot $v$, $w$ or absent; factors $n_{a}$ or $w_{a}v_{a}$ in a slot have weight zero and change nothing). So $Y = 0$ if and only if $r = \tfrac23 p = -(B - L)$, which is (2)'s first statement. A Majorana entry $B(\psi, X\varphi)$ on the sixteen with $\psi = \varphi = \nu^{c}$, the empty state, needs $X$ to raise $B - L$ by the amount $T$ does, and a vev that preserves the Standard Model algebra has $Y = 0$; so the Majorana directions have $Y = 0$ and $B - L = \mp 2$. Then $|p| = 3$, every pair slot carries the same letter, and $r = \pm 2$, both rung slots carrying the same letter; with $Y = 0$ the letters agree, so the class is that of $T$ ($p = 3$, $r = 2$) or of $T^{*}$ ($p = -3$, $r = -2$), up to weight-zero factors. The check finds exactly these two classes among the $243$ and confirms $[B - L, T] = -2T$, $[T_{3R}, T] = T$, $[Y, T] = 0$, $[N_{R}, T] = 2T$ and the opposite weights for $T^{*}$. A single rung vector has $p = 0$, $r = \pm 1$, so $Y = \pm\tfrac12$ and weight $\pm 1$.

(3) The conjugation $*$ is antilinear and reverses products with $v_{a}^{*} = w_{a}$, so it sends weight $q$ to weight $-q$; the scalar part has weight zero, so the scalar part of $X^{*}Y$ vanishes unless $X$ and $Y$ have equal weight, and weight spaces are $*$-orthogonal. For real $s$, $(B_{s}X)^{*} = B_{1/s}(X^{*})$, so for weight $q$, $\|B_{s}X\|_{*} = s^{2q}\|X\|_{*}$. The balanced rung component $c_{H} = H(v_{4} - w_{4})$ (the lift at $\lambda = H$) has components of weights $\pm 1$ with $\|v_{4}\|_{*} = \|w_{4}\|_{*} = \tfrac12$, so $\|B_{s}c_{H}\|_{*} = \tfrac12H^{2}(s^{2} + s^{-2}) = H^{2}\cosh 2\eta = H^{2}\cosh t$, which is $N$ on the lift family (Theorem \claimref{synthesis.rung-unit-is-frame-boost}(4) with $\lambda = He^{\eta}$). For $D = T + T^{*}$, of weights $\pm 2$ with $\|T\|_{*} = \|T^{*}\|_{*} = \tfrac{1}{32}$ (computed exactly), $\|B_{s}D\|_{*}/\|D\|_{*} = \tfrac12(s^{4} + s^{-4}) = \cosh 2t = 2\cosh^{2}t - 1$. For weight $q$ the ratio is $\cosh qt = T_{q}(\cosh t)$. The pairing given by the scalar part of $\widetilde{X}X$ is preserved by every automorphism induced by an isometry; on $c$ it is $Q(c)$, equal to $-H^{2}$ on the lift family, and on $\alpha T + \beta T^{*}$ it is $\alpha\beta$ times the nonzero pairing of $T$ with $T^{*}$, while $T$ alone pairs to zero. The check verifies, for $s = 2, 3, \tfrac32$ and $H = 1$, the two norm formulas, the identity $\|B_{s}D\|_{*}/\|D\|_{*} = 2x^{2} - 1$ with $x = \|B_{s}c_{H}\|_{*}$, its form $\tfrac12 u^{2} - 1$ in $u = 2x = v^{2}/H^{2}$, the vanishing of the cross terms, the invariance of both pairings, and $\|B_{s}T\|_{*} = s^{4}\|T\|_{*}$, which takes every positive value as $s$ varies, so the stored null Majorana element has no size fixed by the frame and lies outside the class of (4).

(4) Under the hypothesis of (4) the Majorana vev is $\kappa D_{H}$ up to the frame and a phase, read in the frame $s = e^{\eta}$ in which the rung component has $N = H^{2}\cosh t$, $v^{2} = 2N$ (the stored normalisation of \claimref{synthesis.rung-field-d-term-action-and-tree-masses-on-the-lift-locus}, in the Hermitian reading of Proposition \claimref{synthesis.rung-field-on-the-lift-locus}; fixed normalisations of $M$ are absorbed in $\kappa$). By (3), $M^{2} = \kappa^{2}H^{2}\cosh 2t = \kappa^{2}H^{2}(2(v^{2}/(2H^{2}))^{2} - 1) = \kappa^{2}(v^{4}/(2H^{2}) - H^{2})$; the term $-\kappa^{2}H^{2}$ is exact and is not a normalisation. For $x = v^{2}/(2H^{2}) \ge 1$, $2x^{2} - 1 \ge x^{2}$, so $M \ge \kappa v^{2}/(2H)$. Record inputs (all stored quantities): $H \le H_{0}$ since $H^{2} = H_{0}^{2}\Omega_{\Lambda}$ with $\Omega_{\Lambda} \le 1$ (the convention of $N$; the two differ by a factor of order one); $\hbar H_{0} < 1.439 \times 10^{-42}$ GeV from the stored Hubble constant, $\hbar$ and megaparsec; $v^{2} = M_{W}^{2}\sin^{2}\theta_{W}\,\alpha^{-1}/\pi$ (the tree relation $M_{W} = gv/2$ with $g = e/\sin\theta_{W}$, cited from memory), above $6.08 \times 10^{4}$ GeV$^{2}$ with $\pi < \tfrac{22}{7}$, which the check proves from Machin's formula by alternating partial sums. Hence, for $\kappa = 1$, $M > 2.1 \times 10^{46}$ GeV, above $1.7 \times 10^{27}$ times the stored Planck mass and above $4.05 \times 10^{34}$ times $4\pi v_{R}$ with the stored $v_{R} = 4.15 \times 10^{10}$ GeV (Proposition \claimref{standard-model.what-angle-asks-neutrino-sector}), the upper end of the admissible range of the Majorana scale in Theorem \claimref{synthesis.majorana-scale-free-on-pati-salam-route}(iii). Under the imported seesaw, $m_{\nu} = -M_{\nu}M_{R}^{-1}M_{\nu}^{\mathsf{T}}$ (Fact \claimref{standard-model.imported-clebsch-structure-seesaw}) with $M_{R} = Mf$, $f$ a family coupling of smallest singular value $\sigma$, and Dirac entries bounded by the perturbative $\sqrt{4\pi}\,v$, every light mass is at most $\|M_{\nu}\|^{2}\|M_{R}^{-1}\| \le 4\pi v^{2}/(\sigma M) \le 8\pi H/\sigma$, and $8\pi H_{0} < m_{3}/(1.38 \times 10^{30})$ with the stored $m_{3} = 0.050$ eV; so the measured $m_{3}$ needs $\sigma < 1/(1.38 \times 10^{30})$. Since $M \ge \kappa v^{2}/(2H)$, bringing $M$ under $4\pi v_{R}$ needs $\kappa < 1/(4.05 \times 10^{34}) < 10^{-34}$. Either rescue introduces a dimensionless number as far from one as $v/H$ is. Neither $N = M_{P}/H$ nor $M_{P}$ enters the relation, which would tie the Majorana scale to $v/H$ alone; since the tie is excluded, the status of the Majorana scale is unaffected by whether $v/H$ is a parameter of the laws or initial data.

(5) If the boost acts on rung slot $4$ only ($v_{4} \mapsto sv_{4}$, $w_{4} \mapsto s^{-1}w_{4}$, all else fixed, an isometry by the argument of (1)), its weight is the slot-$4$ number $n_{4}$, and $T$ has weight $1$; a vev of $T$-type with size fixed in Hubble units in one frame then scales like the rung component, of order $\kappa v$. If instead the boost is extended to the whole algebra as the complexified hypercharge, $\exp(\eta\,\operatorname{ad} Y)$ up to normalisation, every $Y = 0$ element has weight $0$ and is frame-blind, so a tied size stays of order $\kappa H$. A vev along $T$ is the Pati-Salam breaking vev (Proposition \claimref{standard-model.majorana-element}), whose stored value is $v_{R} = 4.15 \times 10^{10}$ GeV; from the stored quantities, with $\pi > 3$ (Machin, partial sums), $v^{2} < 6.37 \times 10^{4}$ GeV$^{2}$, so $v_{R}^{2} > 2.7 \times 10^{16}v^{2}$ ($v_{R}/v > 1.6 \times 10^{8}$) and $v_{R} > 2.8 \times 10^{52}H_{0}$. For $\kappa$ of order one both alternatives are excluded; the check verifies the two weights and these bounds.

Computed, not part of the check: with $H = H_{0}\sqrt{\Omega_{\Lambda}} = 1.19 \times 10^{-42}$ GeV and $v = 246.2$ GeV, $v^{2}/(\sqrt{2}H) = 3.6 \times 10^{46}$ GeV, $35.9$ orders of magnitude above $v_{R}$.